\documentclass[11pt,letterpaper]{article}

\usepackage[utf8]{inputenc}
\usepackage[T1]{fontenc}
\usepackage[spanish,es-tabla,es-nodecimaldot]{babel}
\usepackage[scaled=0.95]{helvet}
\renewcommand{\familydefault}{\sfdefault}
\IfFileExists{beramono.sty}{\usepackage[scaled=0.86]{beramono}}{}
\usepackage[top=2.5cm,bottom=2.5cm,left=2.5cm,right=2.5cm,headheight=15pt]{geometry}
\usepackage[table]{xcolor}
\usepackage{graphicx}
\usepackage{array}
\usepackage{xltabular}
\usepackage{booktabs}
\usepackage{enumitem}
\usepackage{titlesec}
\usepackage{fancyhdr}
\usepackage{lastpage}
\usepackage{tcolorbox}
\usepackage{listings}
\usepackage{tikz}
\usetikzlibrary{positioning,arrows.meta,shapes.geometric,fit,calc}
\PassOptionsToPackage{obeyspaces,spaces}{url}
\usepackage[hidelinks]{hyperref}
\usepackage{parskip}
\usepackage{needspace}

\definecolor{maxamPrimario}{HTML}{870002}
\definecolor{maxamSecundario}{HTML}{D20F12}
\definecolor{maxamTitulo}{HTML}{500203}
\definecolor{maxamTexto}{HTML}{2E2827}
\definecolor{maxamEtiqueta}{HTML}{475569}
\definecolor{maxamFondo}{HTML}{F3F6FF}
\definecolor{maxamBorde}{HTML}{DEE2F0}
\definecolor{maxamCampos}{HTML}{C2C6D4}
\definecolor{codigoFondo}{HTML}{F7F8FC}

\color{maxamTexto}

\titleformat{\section}{\Large\bfseries\color{maxamTitulo}}{\thesection.}{0.6em}{}[{\color{maxamPrimario}\titlerule[0.8pt]}]
\titleformat{\subsection}{\large\bfseries\color{maxamPrimario}}{\thesubsection}{0.6em}{}
\titleformat{\subsubsection}{\normalsize\bfseries\color{maxamTitulo}}{\thesubsubsection}{0.6em}{}
\titlespacing*{\section}{0pt}{18pt}{10pt}
\titlespacing*{\subsection}{0pt}{14pt}{6pt}

\pagestyle{fancy}
\fancyhf{}
\fancyhead[L]{\small\color{maxamEtiqueta}MAXAM FANEXA}
\fancyhead[R]{\small\color{maxamEtiqueta}Sistema de Control de Viáticos --- Documentación General}
\fancyfoot[C]{\small\color{maxamEtiqueta}Página \thepage\ de \pageref{LastPage}}
\renewcommand{\headrulewidth}{0.4pt}
\renewcommand{\headrule}{\hbox to\headwidth{\color{maxamBorde}\leaders\hrule height \headrulewidth\hfill}}

\setlist[itemize]{leftmargin=1.4em,itemsep=2pt,topsep=4pt,label={\color{maxamPrimario}\textbullet}}
\setlist[enumerate]{leftmargin=1.6em,itemsep=2pt,topsep=4pt}

\newcolumntype{L}[1]{>{\raggedright\arraybackslash}p{#1}}
\newcolumntype{Y}{>{\raggedright\arraybackslash}X}
\renewcommand{\arraystretch}{1.35}
\newcommand{\cabecera}{\rowcolor{maxamPrimario}}
\newcommand{\tth}[1]{\textcolor{white}{\textbf{#1}}}
\newcommand{\code}[1]{\nolinkurl{#1}}
\makeatletter
\g@addto@macro{\UrlBreaks}{\do\_\do\-\do\.\do\:\do\=\do\,}
\makeatother

\newtcolorbox{nota}[1][]{colback=maxamFondo,colframe=maxamPrimario,boxrule=0pt,leftrule=3pt,arc=0pt,left=8pt,right=8pt,top=6pt,bottom=6pt,#1}

\lstdefinelanguage{JavaScript}{
  keywords={const,let,var,function,return,if,else,async,await,new,try,catch,finally,throw,for,of,in,true,false,null,undefined,require,module,exports,import,from,export,default},
  sensitive=true,
  comment=[l]{//},
  morecomment=[s]{/*}{*/},
  morestring=[b]',
  morestring=[b]",
  morestring=[b]`
}
\lstset{
  basicstyle=\ttfamily\footnotesize\color{maxamTexto},
  keywordstyle=\color{maxamPrimario}\bfseries,
  commentstyle=\color{maxamEtiqueta}\itshape,
  stringstyle=\color{maxamTitulo},
  backgroundcolor=\color{codigoFondo},
  frame=leftline,
  framerule=2pt,
  rulecolor=\color{maxamPrimario},
  xleftmargin=10pt,
  framexleftmargin=6pt,
  breaklines=true,
  showstringspaces=false,
  columns=fullflexible,
  keepspaces=true,
  aboveskip=8pt,
  belowskip=8pt,
  literate={á}{{\'a}}1 {é}{{\'e}}1 {í}{{\'i}}1 {ó}{{\'o}}1 {ú}{{\'u}}1 {ñ}{{\~n}}1 {Á}{{\'A}}1 {É}{{\'E}}1 {Í}{{\'I}}1 {Ó}{{\'O}}1 {Ú}{{\'U}}1 {Ñ}{{\~N}}1 {—}{{---}}1 {→}{{$\rightarrow$}}1 {├}{{|}}1 {└}{{`}}1 {│}{{|}}1 {─}{{-}}1
}

\hypersetup{pdftitle={Documentación General - Sistema de Control de Viáticos},pdfauthor={Mateo Alejandro Merino Vidal}}

\newcommand{\muestra}[1]{\definecolor{tmpMuestra}{HTML}{#1}\tikz[baseline=-0.6ex]\node[rectangle,fill=tmpMuestra,draw=maxamCampos,minimum width=0.9cm,minimum height=0.38cm]{};}
\newcommand{\estado}[3]{\definecolor{tmpFondo}{HTML}{#1}\definecolor{tmpTexto}{HTML}{#2}\tikz[baseline=-0.6ex]\node[rectangle,rounded corners=5pt,fill=tmpFondo,text=tmpTexto,font=\scriptsize\bfseries,inner sep=3pt]{#3};}
\newcommand{\campo}[3]{\item \code{#1} --- #2 (\texttt{#3}).}
\makeatletter
\renewcommand*\l@subsection{\@dottedtocline{2}{1.5em}{2.6em}}
\renewcommand*\l@subsubsection{\@dottedtocline{3}{4.1em}{4.3em}}
\makeatother

\begin{document}

\begin{titlepage}
\thispagestyle{empty}
\begin{tikzpicture}[remember picture,overlay]
  \fill[maxamPrimario] (current page.north west) rectangle ([yshift=-1.1cm]current page.north east);
  \fill[maxamSecundario] ([yshift=-1.1cm]current page.north west) rectangle ([yshift=-1.3cm]current page.north east);
  \fill[maxamPrimario] (current page.south west) rectangle ([yshift=0.9cm]current page.south east);
\end{tikzpicture}
\centering
\vspace*{1.6cm}
\IfFileExists{logo-maxam.png}{\includegraphics[width=0.55\textwidth]{logo-maxam.png}\par}{{\Huge\bfseries\color{maxamTexto}MAXAM \textcolor{maxamSecundario}{|} FANEXA}\par}
\vspace{1.4cm}
{\color{maxamPrimario}\rule{0.7\textwidth}{1.2pt}}\par
\vspace{0.8cm}
{\fontsize{26}{32}\selectfont\bfseries\color{maxamPrimario}SISTEMA DE CONTROL\\[4pt] DE VIÁTICOS}\par
\vspace{0.5cm}
{\LARGE\bfseries\color{maxamTitulo}Documentación General}\par
\vspace{0.6cm}
{\large\itshape\color{maxamEtiqueta}Gestión de viajes corporativos y rendición de gastos}\par
\vspace{0.8cm}
{\color{maxamPrimario}\rule{0.7\textwidth}{1.2pt}}\par
\vfill
\begin{tabular}{>{\bfseries\color{maxamTitulo}}r L{7.5cm}}
Autor: & Mateo Alejandro Merino Vidal \\
Empresa: & MAXAM FANEXA \\
Versión: & 1.3 \\
Fecha: & Octubre de 2026 \\
\end{tabular}
\vspace{1.8cm}
\end{titlepage}

\tableofcontents
\newpage

\section{Consideraciones}
\begin{enumerate}
  \item El sistema contempla dos tipos de viaje según su alcance geográfico:
  \begin{itemize}
    \item \textbf{Nacional}: los gastos se registran íntegramente en bolivianos.
    \item \textbf{Internacional}: cuenta además con un presupuesto en dólares. Cualquier día se puede gastar en bolivianos o en dólares, y los gastos pagados en otra moneda se anotan como tramos de conversión con el tipo de cambio que le aplicaron al empleado, quien lo escribe en cada tramo.
  \end{itemize}
  \item El medio de transporte se registra por viaje (terrestre, aéreo u otro), junto con la placa cuando se usa un vehículo de la empresa, y se muestra como distintivo en las vistas de revisión.
  \item La rendición de cada viaje reúne:
  \begin{itemize}
    \item El memorándum del viaje (lugar, duración, medio de transporte y objetivo).
    \item Las facturas o recibos de cada gasto, cuando corresponda.
    \item El detalle de los gastos realizados, con su desglose en subgastos si aplica.
  \end{itemize}
  \item El sistema identifica automáticamente los gastos no permitidos, en particular las bebidas alcohólicas, a partir del detalle de productos facturados y de la descripción del gasto. Una rendición con alcohol pasa por una revisión adicional del aprobador.
  \item Al registrar el viaje, el sistema calcula el monto con las tarifas diarias del cargo del empleado; tesorería lo revisa y lo ajusta al asignar los fondos.
  \item El presupuesto se controla día por día y por moneda contra las tarifas diarias del cargo: en un viaje nacional la tarifa en bolivianos; en uno internacional, las dos tarifas por separado. Los hoteles se comparan contra el total asignado. Cada día que excede la tarifa, y el exceso en hoteles, requiere una justificación, y el sistema calcula el saldo a devolver o el importe a reembolsar en cada moneda.
  \item Antes de aprobar o rechazar, los revisores ven la rendición completa: tablas separadas por tipo de gasto, barras de presupuesto, desglose diario y alertas.
  \item Las categorías de gasto incorporan la cuenta contable de Oracle en su nombre, de modo que la planilla exportada resulte directamente utilizable por el área contable.
  \item Se aplican retenciones impositivas automáticas según el tipo de comprobante, conforme a la normativa tributaria boliviana.
  \item Existe un plazo de tolerancia de cuatro días posteriores a la finalización del viaje para completar la carga de gastos. Vencido ese plazo, se debe solicitar una ampliación al revisor.
  \item El empleado puede pedir que otra persona de su sección rinda los gastos de su viaje en su nombre.
  \item El sistema exporta la rendición en Excel, replicando el formato institucional de la planilla de rendición de cuentas, también de varios viajes a la vez para el revisor, el tesorero y el administrador; genera en PDF el memorándum, la confirmación de fondos, los recibos y la rendición final. El administrador puede exportar además a Excel el resumen general, los usuarios, los cargos y las secciones.
\end{enumerate}

\section{Roles del sistema}
El sistema contempla cinco roles y un perfil adicional que se determina por el cargo del usuario.

\begin{xltabular}{\textwidth}{L{1cm} L{3.4cm} Y}
\cabecera \tth{ID} & \tth{Rol} & \tth{Atribuciones} \\
\endfirsthead
\cabecera \tth{ID} & \tth{Rol} & \tth{Atribuciones} \\
\endhead
\rowcolor{maxamFondo} 1 & ADMINISTRADOR & Gestión de usuarios, cargos y secciones; historial de accesos; panel de indicadores filtrable por periodo; exportación a Excel de sus pantallas y de las planillas; días y horas de los recordatorios de pendientes. \\
2 & SUPERVISOR & Revisión previa del viaje y revisión de la rendición de gastos de los empleados a su cargo. \\
\rowcolor{maxamFondo} 3 & EMPLEADO & Creación de viajes, registro de gastos y facturas, solicitudes de ampliación de plazo y de reemplazo. \\
4 & REVISOR & Revisión final de la rendición, resolución de las solicitudes de plazo y de reemplazo y exportación de planillas en grupo. \\
\rowcolor{maxamFondo} 5 & APROBADOR & Aprobación del viaje con emisión del memorándum y revisión adicional de rendiciones con alcohol. \\
--- & TESORERO & Asignación de fondos y exportación de planillas en grupo. Se determina por el cargo (asistente de caja y tesorería), no por el rol. \\
\end{xltabular}

El aprobador y el revisor son roles únicos. El supervisor que revisa cada viaje se determina por jerarquía: primero el jefe directo del empleado; si no lo tiene, los supervisores de su sección; y si tampoco, cualquier supervisor.

\section{Flujo del sistema}
El sistema opera mediante dos flujos de aprobación consecutivos.

\begin{center}
\begin{tikzpicture}[
  est/.style={rectangle,rounded corners=3pt,draw=maxamPrimario,fill=maxamFondo,minimum width=3.4cm,minimum height=0.8cm,font=\scriptsize\bfseries,align=center},
  fin/.style={est,fill=maxamPrimario,text=white},
  rech/.style={est,draw=maxamSecundario,fill=white,text=maxamSecundario},
  fl/.style={-{Stealth[length=2mm]},maxamTitulo,thick},
  node distance=0.35cm]
\node[est] (b) {BORRADOR};
\node[est,below=of b] (rv) {EN\_REVISION\_VIAJE};
\node[est,below=of rv] (av) {APROBADO\_VIAJE};
\node[est,below=of av] (rt) {EN\_REVISION\_TESORERO};
\node[fin,below=of rt] (ec) {EN\_CURSO};
\node[est,right=2.2cm of b] (er) {EN\_REVISION};
\node[est,below=of er] (ea) {EN\_REVISION\_APROBADOR\\(solo con alcohol)};
\node[est,below=of ea] (as) {APROBADO\_SUPERVISOR};
\node[fin,below=of as] (af) {APROBADO\_FINAL};
\node[rech,below=1.1cm of af] (r) {RECHAZADO};
\draw[fl] (b)--(rv); \draw[fl] (rv)--(av); \draw[fl] (av)--(rt); \draw[fl] (rt)--(ec);
\draw[fl] (ec.east) -- ++(1.1,0) |- (er.west);
\draw[fl] (er)--(ea); \draw[fl] (ea)--(as); \draw[fl] (as)--(af);
\draw[fl,dashed,maxamSecundario] (ec.south) |- (r.west);
\node[font=\scriptsize\bfseries,text=maxamPrimario,above=0.25cm of b] {Primer flujo: autorización};
\node[font=\scriptsize\bfseries,text=maxamPrimario,above=0.25cm of er] {Segundo flujo: rendición};
\end{tikzpicture}
\end{center}

\subsection{Primer flujo: autorización del viaje}
\begin{enumerate}
  \item \textbf{Registro del viaje:} el empleado crea el viaje en estado \code{BORRADOR}, indicando motivo, origen, destino, fechas, tipo (nacional o internacional) y medio de transporte.
  \item \textbf{Monto del viaje:} el sistema lo calcula con las tarifas del cargo. En un viaje nacional es días × tarifa en bolivianos; en uno internacional, el día de salida y el de regreso se calculan con la tarifa en bolivianos y los días intermedios con la tarifa en dólares.
  \item \textbf{Envío a revisión:} el viaje pasa a \code{EN_REVISION_VIAJE}.
  \item \textbf{Revisión del supervisor:} el supervisor que corresponde por jerarquía toma el viaje, lo evalúa y lo aprueba o lo rechaza con observaciones. Si lo aprueba, pasa a \code{APROBADO_VIAJE}.
  \item \textbf{Aprobación y memorándum:} el aprobador confirma la solicitud y el sistema emite el memorándum del viaje, que detalla motivo, destino, fechas y transporte autorizado. El viaje pasa a \code{EN_REVISION_TESORERO}.
  \item \textbf{Asignación de fondos:} el tesorero revisa el monto calculado y asigna los montos finales en ambas monedas. Al aprobar, el empleado recibe la confirmación del fondo, el viaje pasa a \code{EN_CURSO} y queda habilitado el registro de gastos.
\end{enumerate}

\subsection{Segundo flujo: rendición de gastos}
\begin{enumerate}
  \item \textbf{Registro de gastos:} durante el viaje, el empleado (o su reemplazo aprobado) registra sus gastos en tres modalidades:
  \begin{itemize}
    \item Gastos sin factura, con cálculo automático de retenciones y emisión de recibo.
    \item Gastos con factura o recibo, cargando la imagen del comprobante y revisando los datos extraídos automáticamente.
    \item Gastos internacionales en dólares, con el registro opcional de tramos de conversión.
  \end{itemize}
  \item \textbf{Validación de fechas:} la fecha de cada gasto debe estar dentro del período del viaje, y la carga debe hacerse dentro del plazo de tolerancia o de la ampliación autorizada.
  \item \textbf{Justificación de excesos:} si algún día supera el viático diario, o los hoteles superan el total asignado, el sistema exige una justificación por cada exceso antes de permitir el envío.
  \item \textbf{Envío a revisión:} el viaje pasa a \code{EN_REVISION}. Las solicitudes de plazo o de reemplazo que hayan quedado sin responder se cierran automáticamente.
  \item \textbf{Revisión del supervisor:} el supervisor toma la rendición y accede a la vista completa, que incluye:
  \begin{itemize}
    \item Datos del empleado y del viaje.
    \item Barras de presupuesto consumido por moneda y desglose por día.
    \item Tarjeta de liquidación con el saldo a devolver o a reembolsar.
    \item Alertas de exceso de presupuesto y de detección de alcohol.
    \item Tablas separadas de gastos con factura, sin factura e internacionales.
    \item Resumen de totales e IVA crédito fiscal.
  \end{itemize}
  Puede observar gastos individuales, exportar la planilla en Excel y aprobar o rechazar. Si aprueba y la rendición tiene alcohol, pasa a \code{EN_REVISION_APROBADOR}; en caso contrario, a \code{APROBADO_SUPERVISOR}.
  \item \textbf{Revisión por alcohol:} el aprobador revisa los gastos con bebidas alcohólicas y aprueba o rechaza. Si aprueba, la rendición pasa a \code{APROBADO_SUPERVISOR}.
  \item \textbf{Revisión final:} el revisor examina la misma información y resuelve la aprobación definitiva, que lleva el viaje a \code{APROBADO_FINAL} y envía al empleado el resultado con el saldo de cada moneda, o su rechazo.
\end{enumerate}

\subsection{Rechazo y ciclos de revisión}
En ambos flujos, el rechazo lleva el viaje al estado \code{RECHAZADO} e incrementa el ciclo de revisión, para que el empleado corrija y reenvíe. Un viaje rechazado antes de iniciarse se corrige editándolo; uno rechazado en la rendición se corrige en sus gastos y vuelve directamente a \code{EN_REVISION}. Las observaciones de ciclos anteriores se conservan como historial, pero solo se consideran las del ciclo vigente: para rechazar, cada revisor debe haber registrado al menos una observación propia en el ciclo actual.

\subsection{Aprobación automática}
Si el responsable de una etapa única (aprobador, tesorero o revisor) es el propio viajero, esa etapa se aprueba sola y los documentos lo indican como aprobación automática. En el caso del tesorero, el fondo queda con el monto calculado al registrar el viaje.

\subsection{Estados del viaje}
\begin{xltabular}{\textwidth}{L{4.6cm} L{3.6cm} Y}
\cabecera \tth{Estado} & \tth{Etiqueta} & \tth{Significado} \\
\endfirsthead
\cabecera \tth{Estado} & \tth{Etiqueta} & \tth{Significado} \\
\endhead
\code{BORRADOR} & \estado{E5E7EB}{374151}{Borrador} & Creado, aún no enviado a revisión. \\
\code{EN_REVISION_VIAJE} & \estado{E8D5FF}{5B00A0}{Revisión Previa} & En revisión previa del supervisor. \\
\code{APROBADO_VIAJE} & \estado{FFE766}{7A5900}{Esperando Aprobador} & Aprobado por el supervisor, pendiente del aprobador. \\
\code{EN_REVISION_TESORERO} & \estado{FFE3C2}{8A4B00}{Esperando Fondos} & Pendiente de asignación de fondos. \\
\code{EN_CURSO} & \estado{D4EDDA}{155724}{En Curso} & Fondos asignados; el empleado registra gastos. \\
\code{EN_REVISION} & \estado{E8D5FF}{5B00A0}{En Revisión} & Rendición enviada, en revisión del supervisor. \\
\code{EN_REVISION_APROBADOR} & \estado{FFE766}{7A5900}{Revisión por Alcohol} & Rendición con alcohol en revisión del aprobador. \\
\code{APROBADO_SUPERVISOR} & \estado{A9C6F6}{000A65}{Apr. Preliminar} & Pendiente de revisión final. \\
\code{APROBADO_FINAL} & \estado{D4EDDA}{155724}{Aprobado} & Rendición aprobada definitivamente. \\
\code{RECHAZADO} & \estado{FFC2C3}{500203}{Rechazado} & Rechazado en cualquiera de los dos flujos. \\
\end{xltabular}

\section{Gestión de plazos}
El plazo de tolerancia concede cuatro días posteriores a la finalización del viaje para completar la carga de gastos. Vencido ese margen, el registro queda bloqueado y se debe pedir una ampliación al revisor, explicando el motivo del retraso. La ampliación vale para el viaje: la puede pedir el titular o su reemplazo, y la respuesta se avisa a ambos. El plazo no corre mientras el viaje está rechazado.

\begin{xltabular}{\textwidth}{L{4.6cm} Y}
\cabecera \tth{Estado de la solicitud} & \tth{Efecto} \\
\endfirsthead
\cabecera \tth{Estado de la solicitud} & \tth{Efecto} \\
\endhead
\rowcolor{maxamFondo} Sin solicitud & Registro bloqueado; se habilita el botón de solicitud. \\
\code{PENDIENTE} & Registro bloqueado hasta la respuesta del revisor. \\
\rowcolor{maxamFondo} \code{APROBADA} vigente & Registro habilitado durante cuatro días desde la respuesta. \\
\code{APROBADA} vencida & Registro bloqueado; puede pedirse otra ampliación. \\
\rowcolor{maxamFondo} \code{RECHAZADA} & Registro bloqueado; puede formularse una nueva solicitud. \\
Cerrada automáticamente & El viaje se envió a revisión antes de la respuesta; ya no aplica. \\
\end{xltabular}

La ampliación extiende el margen para \textit{cargar} gastos, no el rango de fechas admisibles: la fecha de cada gasto debe seguir perteneciendo al período del viaje.

\section{Rendición por terceros}
Mientras el viaje está en curso o rechazado en la rendición, el empleado puede pedir que otra persona activa de su sección rinda los gastos en su nombre. El revisor aprueba o rechaza la solicitud; con la aprobación se avisa al titular y al reemplazo, y el viaje aparece en el panel del reemplazo como ``Rendición de [nombre]''. El titular conserva siempre el acceso a su viaje, y los documentos oficiales llevan su nombre.

\section{Retenciones impositivas}
El sistema calcula automáticamente las retenciones según el tipo de comprobante, conforme a la normativa boliviana. Los valores se guardan en el momento del registro para que los reportes históricos no cambien si cambian las alícuotas.

\begin{xltabular}{\textwidth}{L{1.6cm} L{5.4cm} Y}
\cabecera \tth{Código} & \tth{Tipo de comprobante} & \tth{Retenciones aplicadas} \\
\endfirsthead
\cabecera \tth{Código} & \tth{Tipo de comprobante} & \tth{Retenciones aplicadas} \\
\endhead
\rowcolor{maxamFondo} \code{F} & Compra de bien o servicio con factura & Ninguna; genera IVA crédito fiscal (13\% sobre el subtotal) \\
\code{R} & Recibo: documento sin IVA ni retención & Ninguna; no lleva IVA \\
\rowcolor{maxamFondo} \code{C} & Compra de bien sin factura & IUE 5\% e IT 3\% sobre el total \\
\code{S} & Servicio o alquiler sin factura & RC-IVA 13\% e IT 3\% sobre el total \\
\end{xltabular}

Los gastos internacionales quedan exentos de retenciones. Los gastos con alcohol pierden toda retención y crédito fiscal y se imputan íntegros como costo. En las facturas (\code{F}), el total es el subtotal más el IVA del 13\%; al empleado se le reconoce el total, y en la planilla el IVA figura en su columna como crédito fiscal de la empresa y el importe costo/gasto es el total menos el IVA. Los recibos (\code{R}) no llevan IVA: el monto es el total pagado y la planilla muestra IVA cero. En los gastos sin factura, las retenciones se calculan sobre el total del documento y al proveedor se le cancela el total menos las retenciones. Por ejemplo, en un servicio de 47.62 el RC-IVA es 6.19, el IT 1.43 y el importe a cancelar 40.00. El total cuenta para el presupuesto y como costo en la planilla.

\section{Categorías de gasto}
Las categorías incorporan en su nombre el código de la cuenta contable de Oracle, lo que permite que la planilla exportada sea directamente utilizable por el área contable. Al agregar nuevas categorías debe conservarse ese formato.

\begin{xltabular}{\textwidth}{Y}
\cabecera \tth{Categoría} \\
\endfirsthead
\cabecera \tth{Categoría} \\
\endhead
\rowcolor{maxamFondo} 626000 BILLETES (VIAJE) \\
626010 TAXIS (VIAJE) --- no requiere comprobante \\
\rowcolor{maxamFondo} 626020 COCHE PROPIO-KMS (VIAJE) \\
626030 HOTELES (VIAJE) \\
\rowcolor{maxamFondo} 626040 MANUTENCION (VIAJE) \\
626050 VEHICULOS DE ALQUILER (VIAJE) \\
\rowcolor{maxamFondo} 626060 PEAJE AUTOPISTAS (VIAJE) \\
626070 APARCAMIENTO Y OTROS GASTOS (VIAJE) \\
\rowcolor{maxamFondo} 626100 REUN.SINDIC.MANUTENCION/AT.TERCEROS (VIAJE) \\
OTROS \\
\end{xltabular}

\section{Notificaciones}
\begin{itemize}
  \item \textbf{Avisos inmediatos:} el empleado recibe el memorándum, la confirmación del fondo, el resultado de su rendición, los recibos y la respuesta a sus solicitudes; el revisor recibe las nuevas solicitudes de plazo y de reemplazo; el reemplazo recibe el aviso de la rendición asignada.
  \item \textbf{Resumen de pendientes:} en los días y horas que el administrador configura en la pantalla \emph{Recordatorios} (por defecto, de lunes a viernes a las 08:00, 12:00 y 16:00, hora de Bolivia), cada supervisor, aprobador, revisor y tesorero con trabajo pendiente recibe un solo correo con una línea por tipo de pendiente.
\end{itemize}

\section{Arquitectura del sistema}

\subsection{Modelo C4}
La arquitectura del sistema se documentó siguiendo el modelo C4, que describe el software en niveles sucesivos de detalle: contexto del sistema, contenedores, componentes y código. Los diagramas pueden consultarse en el siguiente enlace:

\begin{center}
\url{https://drive.google.com/file/d/1nhu-JCIgQiKUj8bUSOvvz443vKSRAN_N/view?usp=sharing}
\end{center}

\subsection{Organización del frontend}
El cliente web combina una organización por dominios con una jerarquía de componentes basada en Atomic Design. La presentación está separada de la lógica: los componentes visuales no llaman a la API; esa responsabilidad recae en hooks que consumen la capa de servicios.

\begin{xltabular}{\textwidth}{L{3.2cm} Y}
\cabecera \tth{Capa} & \tth{Responsabilidad} \\
\endfirsthead
\cabecera \tth{Capa} & \tth{Responsabilidad} \\
\endhead
\rowcolor{maxamFondo} Página & Invoca el hook y distribuye sus valores entre los componentes. \\
Hook & Mantiene el estado, valida y llama a los servicios. \\
\rowcolor{maxamFondo} Servicio & Encapsula el endpoint de la API y normaliza la respuesta. \\
Componente & Dibuja la interfaz a partir de las propiedades recibidas. \\
\end{xltabular}

\subsection{Organización del backend}
El servidor sigue una arquitectura en capas: la ruta declara el endpoint y encadena los middlewares; los middlewares validan la sesión, el rol y el estado del usuario; el controlador lee los parámetros y delega; y el servicio aplica las reglas de negocio y accede a los datos.

\subsection{Seguridad}
\begin{itemize}
  \item La sesión usa un token de acceso de 15 minutos guardado en memoria y un token de renovación de 7 días en una cookie \code{httpOnly}. Cambiar el rol, suspender una cuenta o cambiar la contraseña cierra las demás sesiones al instante.
  \item Las contraseñas temporales (cuenta nueva, clave puesta por el administrador o ingreso con código de recuperación) deben cambiarse antes de usar el sistema, y toda contraseña vence a los 90 días.
  \item Los intentos de ingreso, de recuperación y de lectura de facturas tienen un límite por periodo. El de ingreso es por cuenta: si un usuario se equivoca muchas veces, solo su cuenta espera 15 minutos y los demás siguen entrando.
  \item Cada ingreso, salida y cambio de contraseña queda registrado en el historial de accesos.
\end{itemize}

\section{Diseño de interfaz}
El flujo de pantallas se diseñó en Figma buscando una experiencia clara, moderna y consistente entre dispositivos, priorizando la legibilidad, la jerarquía visual y la coherencia en el uso de tipografías y colores.

\begin{nota}
\textbf{Alcance del prototipo.} El archivo de Figma contiene únicamente un \emph{prototipo inicial}, elaborado durante la etapa de concepción. No representa todas las pantallas ni el comportamiento definitivo de la aplicación: el sistema final incorpora muchas vistas, flujos y componentes que no figuran en el prototipo, y varios de los que sí aparecen se rediseñaron durante la implementación a medida que evolucionaron los requerimientos. Su valor documental es el de una guía visual de la identidad gráfica adoptada, no el de una especificación de la interfaz construida.
\end{nota}

\begin{center}
\url{https://www.figma.com/design/2t3LLwLBtfypltzjqcqSGS/Pantallas_MAXAM_IT?node-id=0-1&t=P0O5UMwvyqVZ0A7N-1}
\end{center}

\subsection{Tipografías}
\begin{itemize}
  \item \textbf{Inter}: títulos, subtítulos, etiquetas y contenido general, con una apariencia profesional, limpia y muy legible.
  \item \textbf{Nunito}: botones y elementos de acción, como contraste amigable frente a la tipografía del contenido.
\end{itemize}

\subsection{Tamaños de fuente}
\begin{xltabular}{\textwidth}{L{5cm} L{4cm} Y}
\cabecera \tth{Elemento} & \tth{Escritorio} & \tth{Móvil} \\
\endfirsthead
\cabecera \tth{Elemento} & \tth{Escritorio} & \tth{Móvil} \\
\endhead
\rowcolor{maxamFondo} Títulos & 45 px Bold & 36 px Bold \\
Subtítulos & 36 px SemiBold & 24 px SemiBold \\
\rowcolor{maxamFondo} Texto destacado & --- & 30 px \\
h1 (texto importante) & 24 px & 20 px \\
\rowcolor{maxamFondo} h2 (texto normal) & 20 px & 18 px \\
h3 (subtexto) & 16 px & 14 px \\
\rowcolor{maxamFondo} h4 (texto mínimo) & 14 px & 12 px \\
Botones & --- & 16 px \\
\end{xltabular}

\subsection{Paleta de colores}
La paleta se centraliza en un único archivo de constantes del cual todos los componentes obtienen sus valores. Ningún componente define colores sueltos, lo que garantiza la consistencia y permite modificar la identidad visual desde un solo punto.

\begin{xltabular}{\textwidth}{L{1.3cm} L{4cm} L{2cm} Y}
\cabecera \tth{Color} & \tth{Token} & \tth{Valor} & \tth{Uso} \\
\endfirsthead
\cabecera \tth{Color} & \tth{Token} & \tth{Valor} & \tth{Uso} \\
\endhead
\muestra{870002} & \code{primary} & \texttt{\#870002} & Color institucional; acciones principales y ventanas emergentes. \\
\muestra{D20F12} & \code{secondary} & \texttt{\#D20F12} & Botón de confirmación, acciones destructivas y errores. \\
\muestra{500203} & \code{title} & \texttt{\#500203} & Títulos y encabezados de sección. \\
\muestra{2E2827} & \code{text} & \texttt{\#2E2827} & Texto principal. \\
\muestra{475569} & \code{labels} & \texttt{\#475569} & Texto secundario y etiquetas de campo. \\
\muestra{FFFFFF} & \code{background} & \texttt{\#FFFFFF} & Fondo general. \\
\muestra{000000} & \code{backgroundSecondary} & \texttt{\#000000} & Íconos sobre círculo blanco y cabeceras de alto contraste. \\
\muestra{F3F6FF} & \code{backgroundHeader} & \texttt{\#F3F6FF} & Tarjetas y encabezados de tabla. \\
\muestra{DEE2F0} & \code{dataFields} & \texttt{\#DEE2F0} & Bordes y separadores. \\
\muestra{C2C6D4} & \code{fields} & \texttt{\#C2C6D4} & Controles deshabilitados. \\
\muestra{FDE9E9} & \code{error} & \texttt{\#FDE9E9} & Fondo de mensajes de error. \\
\muestra{F2D9C5} & \code{travelTypes} & \texttt{\#F2D9C5} & Fondo del distintivo de tipo de viaje. \\
\muestra{A45718} & \code{travelTypesText} & \texttt{\#A45718} & Texto del distintivo de tipo de viaje. \\
\muestra{A69896} & \code{environmentTypes} & \texttt{\#A69896} & Fondo del distintivo de transporte. \\
\muestra{47312E} & \code{environmentTypesText} & \texttt{\#47312E} & Texto del distintivo de transporte. \\
\muestra{F1F2FF} & \code{positionRole} & \texttt{\#F1F2FF} & Fondo de distintivos de cargo y rol. \\
\muestra{E9EBF2} & \code{footer} & \texttt{\#E9EBF2} & Fondo del pie de página. \\
\muestra{DCDCDC} & \code{bar} & \texttt{\#DCDCDC} & Fondo de las barras de progreso. \\
\end{xltabular}

\subsection{Colores de estado}
Cada estado tiene un único nombre y un único par de colores en todas las pantallas y roles. Los avisos de confirmación o rechazo usan el color del estado en que queda el viaje.

\begin{xltabular}{\textwidth}{L{4.2cm} L{3.8cm} L{2.4cm} Y}
\cabecera \tth{Grupo} & \tth{Muestra} & \tth{Fondo} & \tth{Texto} \\
\endfirsthead
\cabecera \tth{Grupo} & \tth{Muestra} & \tth{Fondo} & \tth{Texto} \\
\endhead
Borrador & \estado{E5E7EB}{374151}{Borrador} & \texttt{\#E5E7EB} & \texttt{\#374151} \\
Pendiente de revisión & \estado{E8D5FF}{5B00A0}{En Revisión} & \texttt{\#E8D5FF} & \texttt{\#5B00A0} \\
Esperando aprobador & \estado{FFE766}{7A5900}{Esperando Aprobador} & \texttt{\#FFD700AA} & \texttt{\#7A5900} \\
Esperando fondos & \estado{FFE3C2}{8A4B00}{Esperando Fondos} & \texttt{\#FFD8A8AA} & \texttt{\#8A4B00} \\
Pendiente de revisión final & \estado{A9C6F6}{000A65}{Apr. Preliminar} & \texttt{\#85AFF3AB} & \texttt{\#000A65} \\
En curso y aprobado & \estado{D4EDDA}{155724}{Aprobado} & \texttt{\#D4EDDA} & \texttt{\#155724} \\
Rechazado & \estado{FFC2C3}{500203}{Rechazado} & \texttt{\#FFA7A8AA} & \texttt{\#500203} \\
Observado & \estado{FEF3CD}{856404}{Observado} & \texttt{\#FEF3CD} & \texttt{\#856404} \\
\end{xltabular}

En las pantallas del empleado, ``En Curso'' usa el rojo institucional porque marca su viaje activo.

\subsection{Uso del color}
\begin{itemize}
  \item Los \textbf{fondos claros} predominan para favorecer la lectura del contenido.
  \item El \textbf{rojo institucional} identifica las acciones principales, la tarjeta del viaje en curso y todas las ventanas emergentes, con el ícono en un círculo blanco.
  \item El \textbf{rojo secundario} se reserva para la confirmación dentro de las ventanas, las acciones destructivas y los errores.
  \item El \textbf{color de títulos} refuerza la jerarquía tipográfica sobre fondos neutros.
  \item Los \textbf{colores neutros} ordenan campos, divisores y zonas de datos sin generar ruido visual.
  \item Los \textbf{colores de estado} muestran de inmediato en qué etapa del flujo está cada viaje.
\end{itemize}

\subsection{Criterios generales de diseño}
\begin{itemize}
  \item \textbf{Jerarquía visual clara}: tamaños y pesos tipográficos consistentes para guiar al usuario.
  \item \textbf{Consistencia}: los mismos colores, tipografías, contadores (``Mostrando X de Y'') y ventanas en todas las pantallas.
  \item \textbf{Adaptabilidad}: interfaz adaptable a escritorio y celular.
  \item \textbf{Legibilidad}: tipografías y contrastes elegidos para facilitar la lectura.
  \item \textbf{Experiencia de usuario}: navegación intuitiva, actualización automática de las bandejas cada 30 segundos, respuesta inmediata a cada acción y transiciones animadas entre pantallas y ventanas.
\end{itemize}

\section{Planillas}
Información confidencial, contenida en la documentación de Excel recibida. El formato definitivo de la planilla lo proporcionará la empresa y se incorporará cuando se reciba. El sistema replica el formato institucional al exportar la rendición, incorporando el logotipo, la tabla de simbología de comprobantes, la distinción por color entre gastos nacionales e internacionales, las listas desplegables de cuenta contable, las justificaciones de exceso y los bloques de firma de cada instancia de aprobación. El empleado descarga además su planilla en PDF una vez aprobado el viaje.

\section{Stack tecnológico}

\subsection{Frontend}
\begin{itemize}
  \item \textbf{React}: biblioteca base de la interfaz.
  \item \textbf{Vite}: compilación y servidor de desarrollo.
  \item \textbf{Tailwind CSS}: utilidades de maquetado.
  \item \textbf{React Router DOM}: enrutamiento y protección de rutas por rol.
  \item \textbf{Framer Motion}: animaciones de pantallas y ventanas.
  \item \textbf{Axios}: cliente HTTP con renovación automática de la sesión.
  \item \textbf{Recharts}: gráficos del panel administrativo.
  \item \textbf{ExcelJS}: generación de la planilla de rendición.
  \item \textbf{Lucide React}: iconografía.
  \item \textbf{Storybook}: documentación de componentes.
  \item \textbf{Vitest} y \textbf{Playwright}: pruebas unitarias y de extremo a extremo.
\end{itemize}

\subsection{Backend}
\begin{itemize}
  \item \textbf{Node.js} y \textbf{Express}: entorno de ejecución y servidor de la API REST.
  \item \textbf{jsonwebtoken} y \textbf{bcrypt}: autenticación y cifrado de contraseñas.
  \item \textbf{Multer}: recepción de imágenes de comprobantes y fotos de perfil.
  \item \textbf{Puppeteer}: generación de documentos PDF.
  \item \textbf{node-cron}: envío programado del resumen de pendientes.
  \item \textbf{Helmet}, \textbf{CORS} y \textbf{express-rate-limit}: cabeceras de seguridad, orígenes permitidos y límite de intentos.
  \item \textbf{Jest} y \textbf{Supertest}: pruebas unitarias y de la API.
\end{itemize}

\subsection{Servicios externos}
\begin{itemize}
  \item \textbf{Supabase}: base de datos PostgreSQL y almacenamiento de archivos.
  \item \textbf{Google Gemini}: extracción de datos de comprobantes y detección de bebidas alcohólicas.
  \item \textbf{Brevo}: envío de correo transaccional.
  \item \textbf{AbstractAPI}: verificación de correos antes de enviar códigos de recuperación.
\end{itemize}

\subsection{Herramientas}
\begin{itemize}
  \item \textbf{GitHub}: control de versiones; \textbf{GitHub Actions} ejecuta las pruebas automáticas en cada envío.
  \item \textbf{Figma}: prototipado de la interfaz.
  \item \textbf{dbdiagram.io}: modelado del esquema relacional.
  \item \textbf{\LaTeX}: redacción de la documentación.
  \item \textbf{Vercel}: despliegue del cliente web.
  \item \textbf{Render}: despliegue del servidor durante la prueba piloto.
\end{itemize}

\section{Pruebas automáticas}
El sistema cuenta con tres niveles de pruebas automáticas, que verifican su funcionamiento sin intervención manual:
\begin{itemize}
  \item \textbf{Unitarias}: comprueban reglas puntuales, como las retenciones, el control diario por moneda, el IVA de facturas y recibos y las cantidades con decimales del detalle (Jest en el servidor y Vitest en el cliente).
  \item \textbf{De la API}: ejecutan los servicios del servidor contra una base de prueba y comprueban permisos, montos y respuestas (Jest y Supertest).
  \item \textbf{De extremo a extremo}: recorren el sistema en un navegador real, en escritorio y celular, como lo haría cada rol (Playwright).
\end{itemize}
Las pruebas se ejecutan automáticamente en GitHub Actions cada vez que se envían cambios al repositorio, sobre una base de datos temporal que se crea y se elimina en cada ejecución, sin tocar la base real.

\section{Extracción de datos de comprobantes}
Al cargar una factura, la imagen se envía al modelo de visión de Google Gemini con una instrucción que fija el contexto tributario boliviano: el término ``Importe'' corresponde al total con impuesto incluido y la alícuota del IVA es del trece por ciento. El modelo devuelve proveedor, NIT, número de factura, fecha de emisión, monto sin impuestos, IVA, total, tipo de documento y detalle de productos.

\begin{enumerate}
  \item La llamada reintenta ante errores temporales y, si el modelo principal falla o agota su cuota, usa un modelo de respaldo.
  \item Los datos que faltan se marcan como ``No Especificado'', la fecha se valida y la cantidad de cada producto se redondea a tres decimales. En un recibo, el monto es el total pagado y no se descuenta IVA. Si un recibo o boleto no tiene número, se registra su código, que puede llevar símbolos como @ o \$.
  \item El empleado revisa y corrige los datos en una vista de tres columnas (imagen, formulario y productos) antes de guardar, incluida la fecha de emisión por si se leyó mal. Al modificar después una factura ya guardada, la fecha queda bloqueada.
  \item La cantidad de cada producto admite hasta tres decimales, por ejemplo 53.345 litros de combustible o 6 unidades.
  \item En el formulario, tanto al escanear como al cargar una factura a mano, se registra el total que figura en el papel: en una factura el sistema separa por dentro el subtotal y el IVA sin cambiar el total; en un recibo no hay IVA.
  \item Una factura repetida (mismo número, fecha y proveedor) se rechaza.
\end{enumerate}

\section{Despliegue}
La aplicación se distribuye en dos servicios independientes, el cliente web y el servidor de aplicaciones, cada uno en una plataforma especializada. La base de datos y el almacenamiento de archivos permanecen en Supabase.

\subsection{Cliente web}
El cliente se despliega en \textbf{Vercel}, plataforma orientada a aplicaciones de frontend que compila automáticamente a partir del repositorio y distribuye el sitio mediante una red de entrega de contenidos. Cada cambio incorporado a la rama configurada desencadena una nueva compilación y publicación.

\begin{center}
\url{https://scv-frontend-git-develop-mat13.vercel.app}
\end{center}

\subsection{Servidor de aplicaciones}
Durante la prueba piloto, el servidor se despliega en \textbf{Render}, que ejecuta el proceso de Node.js y lo reinicia ante fallos o nuevas publicaciones. Para el uso definitivo se prevé un servidor VPS, donde el proceso queda siempre activo; el procedimiento está detallado en el manual técnico del backend.

\begin{center}
\url{https://scv-backend-7yft.onrender.com}
\end{center}

\subsection{Variables de entorno}
Ninguna credencial se incluye en el repositorio. Ambos servicios obtienen su configuración de las variables de entorno definidas en cada plataforma.

\begin{xltabular}{\textwidth}{L{5.2cm} Y L{2cm}}
\cabecera \tth{Variable} & \tth{Propósito} & \tth{Servicio} \\
\endfirsthead
\cabecera \tth{Variable} & \tth{Propósito} & \tth{Servicio} \\
\endhead
\rowcolor{maxamFondo} \code{SUPABASE_URL} & Dirección del proyecto de base de datos. & Backend \\
\code{SUPABASE_KEY} & Clave de servicio para el acceso a los datos. & Backend \\
\rowcolor{maxamFondo} \code{JWT_SECRET} & Firma del token de acceso. & Backend \\
\code{JWT_REFRESH_SECRET} & Firma del token de renovación de sesión. & Backend \\
\rowcolor{maxamFondo} \code{GEMINI_API_KEY} & Extracción de comprobantes y detección de alcohol. & Backend \\
\code{BREVO_API_KEY} & Envío de correo transaccional. & Backend \\
\rowcolor{maxamFondo} \code{ABSTRACT_EMAIL_API_KEY} & Verificación de direcciones de correo. & Backend \\
\code{FRONTEND_URL} & Enlace del botón ``Ingresar al sistema'' de los correos. & Backend \\
\rowcolor{maxamFondo} \code{NODE_ENV} & \code{production} activa la cookie segura de sesión. & Backend \\
\code{PORT} & Puerto del servidor. & Backend \\
\rowcolor{maxamFondo} \code{CORS_ORIGINS} & Opcional: dominios del cliente permitidos. & Backend \\
\code{TRUST_PROXY} & Opcional: servidor detrás de un proxy (VPS). & Backend \\
\rowcolor{maxamFondo} \code{VITE_API_URL} & Dirección del servidor de aplicaciones. & Frontend \\
\end{xltabular}

\subsection{Consideraciones de despliegue}
\begin{itemize}
  \item \textbf{Orígenes permitidos}: el servidor solo acepta peticiones de los dominios del cliente declarados en su configuración. Si el cliente cambia de dominio, debe agregarse en \code{CORS_ORIGINS}; de lo contrario, el navegador rechazará las respuestas.
  \item \textbf{Enrutamiento del cliente}: al ser una aplicación de una sola página, la plataforma debe redirigir todas las rutas al documento principal para que funcionen los enlaces directos a una pantalla interna.
  \item \textbf{HTTPS}: la cookie de sesión exige que el servidor se publique con HTTPS cuando el cliente está en otro dominio.
  \item \textbf{Latencia inicial}: en los planes gratuitos, el servidor puede suspenderse tras un período sin uso, lo que demora la primera petición y puede retrasar el resumen de pendientes. En un VPS esto no ocurre.
\end{itemize}

\section{Base de datos}

\subsection{Herramienta de diseño}
El esquema relacional se modeló con \textbf{dbdiagram.io}, herramienta web para construir diagramas entidad-relación a partir de definiciones en DBML, con importación directa de sentencias SQL. El diagrama del proyecto puede consultarse en:

\begin{center}
\url{https://dbdiagram.io/d/SCV-DATABASE_DIAGRAM-6aa21bcd70fd27e3c764a262}
\end{center}

\subsection{Motor de base de datos}
El sistema utiliza \textbf{PostgreSQL}. Se aprovechan características propias del motor: \texttt{SERIAL} para generar identificadores, \texttt{DECIMAL(10,2)} y \texttt{NUMERIC} para valores monetarios con precisión fija, restricciones \texttt{CHECK} para valores enumerados y \texttt{TIMESTAMPTZ} para fechas y horas.

El uso de \texttt{TIMESTAMPTZ} en lugar de \texttt{TIMESTAMP} resultó determinante: con el tipo sin zona horaria, PostgreSQL descartaba el huso al guardar el valor y el cliente lo interpretaba después como hora local, desplazando las fechas cuatro horas respecto del horario de Bolivia.

\subsection{Tablas del esquema}
El esquema está compuesto por \textbf{24 tablas}, agrupadas según las entidades principales del dominio.

\begin{xltabular}{\textwidth}{L{3.2cm} Y}
\cabecera \tth{Agrupación} & \tth{Tablas} \\
\endfirsthead
\cabecera \tth{Agrupación} & \tth{Tablas} \\
\endhead
\rowcolor{maxamFondo} Catálogo & Rol, Cargo, Seccion, Categoria\_Gasto, Impuesto, Proveedor \\
Usuarios & Usuario, Codigo\_Verificacion, Auditoria \\
\rowcolor{maxamFondo} Viajes & Viaje, Solicitud\_Autorizacion\_Plazo, Solicitud\_Reemplazo, Revision\_Viaje \\
Gastos & Gasto, Gasto\_Subitem, Gasto\_Tramo\_Moneda, Imagen \\
\rowcolor{maxamFondo} Facturas & Factura, Detalle\_Factura, Factura\_Impuestos \\
Transversal & Comentario, Correlativo\_Recibo, Recibo, Configuracion\_Recordatorio \\
\end{xltabular}

\subsubsection{Rol}
Roles del sistema, que determinan los permisos de cada usuario. Sus identificadores son fijos y no deben alterarse.
\begin{itemize}
  \campo{id_rol}{Identificador único}{SERIAL, clave primaria}
  \campo{nombre}{Nombre del rol}{VARCHAR(20), único}
\end{itemize}

\subsubsection{Cargo}
Puestos laborales. Cada cargo tiene un viático diario en ambas monedas, usado para el presupuesto sugerido y el control diario.
\begin{itemize}
  \campo{id_cargo}{Identificador único}{SERIAL, clave primaria}
  \campo{nombre}{Nombre del cargo}{VARCHAR(100), único}
  \campo{monto_diario}{Viático diario en bolivianos}{DECIMAL(10,2)}
  \campo{monto_diario_usd}{Viático diario en dólares}{NUMERIC(10,2)}
  \campo{activo}{Indica si el cargo está habilitado}{BOOLEAN}
\end{itemize}

\subsubsection{Seccion}
Secciones de la empresa. Se usan como respaldo en la asignación del supervisor y para elegir el reemplazo.
\begin{itemize}
  \campo{id_seccion}{Identificador único}{SERIAL, clave primaria}
  \campo{nombre}{Nombre de la sección}{VARCHAR(100), único}
  \campo{activo}{Indica si la sección está habilitada}{BOOLEAN}
\end{itemize}

\subsubsection{Usuario}
Información de todos los usuarios registrados.
\begin{itemize}
  \campo{id_usuario}{Identificador único}{SERIAL, clave primaria}
  \campo{email_corporativo}{Correo institucional}{VARCHAR(100), único}
  \campo{nombre}{Nombre}{VARCHAR(30)}
  \campo{apellido_paterno}{Primer apellido}{VARCHAR(30)}
  \campo{apellido_materno}{Segundo apellido, opcional}{VARCHAR(30)}
  \campo{telefono}{Teléfono de 7 u 8 dígitos, opcional}{VARCHAR(8)}
  \campo{carnet_identidad}{Cédula de identidad}{VARCHAR(15)}
  \campo{activo}{Indica si el usuario está habilitado}{BOOLEAN}
  \campo{contrasenia}{Hash de la contraseña}{TEXT}
  \campo{foto_perfil}{URL de la fotografía, opcional}{TEXT}
  \campo{id_jefe_directo}{Jefe directo, para la asignación del supervisor}{INTEGER}
  \campo{id_seccion}{Sección del usuario}{INTEGER}
  \campo{ultima_cambio_contrasenia}{Último cambio de clave}{TIMESTAMPTZ}
  \campo{refresh_token_invalido_desde}{Invalida las sesiones anteriores a esta marca}{TIMESTAMPTZ}
  \campo{debe_cambiar_contrasenia}{Contraseña temporal pendiente de cambio}{BOOLEAN}
  \campo{motivo_cambio_contrasenia}{TEMPORAL o RECUPERACION}{VARCHAR(15)}
  \campo{id_cargo}{Cargo del usuario}{INTEGER}
  \campo{id_rol}{Rol del usuario}{INTEGER}
\end{itemize}

\subsubsection{Viaje}
Entidad central del sistema. Su estado gobierna ambos flujos de aprobación.
\begin{itemize}
  \campo{id_viaje}{Identificador único}{SERIAL, clave primaria}
  \campo{motivo}{Propósito del viaje, hasta 500 caracteres}{VARCHAR(500)}
  \campo{origen}{Localidad de partida, opcional}{TEXT}
  \campo{destino}{Localidad de destino}{VARCHAR(200)}
  \campo{fecha_inicio}{Fecha de inicio}{DATE}
  \campo{fecha_fin}{Fecha de finalización}{DATE}
  \campo{tipo}{Nacional o Internacional}{VARCHAR(15), CHECK}
  \campo{transporte}{Medio de transporte}{TEXT, Terrestre por defecto}
  \campo{placa_vehiculo}{Placa del vehículo de la empresa, opcional}{VARCHAR(20)}
  \campo{monto_asignado}{Presupuesto en bolivianos, calculado con las tarifas del cargo y ajustado por tesorería}{DECIMAL(10,2)}
  \campo{monto_asignado_usd}{Presupuesto en dólares, calculado con las tarifas del cargo y ajustado por tesorería}{NUMERIC(10,2)}
  \campo{estado}{Uno de los diez estados definidos}{VARCHAR(25), CHECK}
  \campo{fue_iniciado}{Distingue un rechazo en la autorización de uno en la rendición}{BOOLEAN}
  \campo{ciclo_revision}{Contador de rondas de revisión}{INTEGER}
  \campo{tiene_alcohol}{Resultado del análisis de los gastos}{BOOLEAN}
  \campo{id_usuario}{Empleado titular}{INTEGER}
  \campo{id_supervisor_asignado}{Supervisor que tomó la revisión}{INTEGER}
  \campo{id_aprobador_asignado}{Aprobador asignado}{INTEGER}
  \campo{id_revisor_asignado}{Revisor de la instancia final}{INTEGER}
  \campo{id_tesorero_asignado}{Tesorero que asignó los fondos}{INTEGER}
\end{itemize}

\subsubsection{Solicitud\_Autorizacion\_Plazo}
Solicitudes de ampliación del plazo de carga de gastos.
\begin{itemize}
  \campo{id_solicitud}{Identificador único}{SERIAL, clave primaria}
  \campo{id_viaje}{Viaje asociado}{INTEGER, borrado en cascada}
  \campo{id_empleado}{Quien la pidió: titular o reemplazo}{INTEGER}
  \campo{motivo}{Justificación del retraso}{VARCHAR(500)}
  \campo{estado}{PENDIENTE, APROBADA o RECHAZADA}{VARCHAR(20)}
  \campo{id_revisor}{Revisor que la resolvió; vacío si se cerró automáticamente}{INTEGER}
  \campo{observacion_revisor}{Motivo del rechazo o del cierre}{VARCHAR(500)}
  \campo{fecha_solicitud}{Fecha de la solicitud}{TIMESTAMPTZ}
  \campo{fecha_respuesta}{Fecha de la respuesta}{TIMESTAMPTZ}
\end{itemize}

\subsubsection{Solicitud\_Reemplazo}
Solicitudes para que otra persona rinda los gastos de un viaje.
\begin{itemize}
  \campo{id_solicitud}{Identificador único}{SERIAL, clave primaria}
  \campo{id_viaje}{Viaje asociado}{INTEGER, borrado en cascada}
  \campo{id_solicitante}{Titular del viaje}{INTEGER}
  \campo{id_sustituto}{Persona propuesta como reemplazo}{INTEGER}
  \campo{estado}{PENDIENTE, APROBADA o RECHAZADA}{VARCHAR(20)}
  \campo{id_revisor}{Revisor que la resolvió; vacío si se cerró automáticamente}{INTEGER}
  \campo{observacion_revisor}{Motivo del rechazo o del cierre}{VARCHAR(500)}
  \campo{fecha_solicitud}{Fecha de la solicitud}{TIMESTAMPTZ}
  \campo{fecha_respuesta}{Fecha de la respuesta}{TIMESTAMPTZ}
\end{itemize}

\subsubsection{Revision\_Viaje}
Historial de aprobaciones y rechazos, usado en las pestañas ``Aprobados'' y ``Rechazados'' de cada revisor.
\begin{itemize}
  \campo{id_revision}{Identificador único}{SERIAL, clave primaria}
  \campo{id_viaje}{Viaje revisado}{INTEGER, borrado en cascada}
  \campo{id_usuario}{Persona que revisó}{INTEGER}
  \campo{etapa}{Etapa del flujo en que se revisó}{VARCHAR(20), CHECK}
  \campo{accion}{APROBADO o RECHAZADO}{VARCHAR(10)}
  \campo{automatica}{Etapa aprobada sola porque el responsable era el viajero}{BOOLEAN}
  \campo{fecha}{Fecha de la revisión}{TIMESTAMPTZ}
\end{itemize}

\subsubsection{Categoria\_Gasto}
Clasifica los gastos por concepto, con la cuenta contable de Oracle en el nombre.
\begin{itemize}
  \campo{id_categoria}{Identificador único}{SERIAL, clave primaria}
  \campo{nombre}{Denominación con código contable}{VARCHAR(50), único}
  \campo{requiere_comprobante}{Indica si la categoría exige comprobante}{BOOLEAN}
\end{itemize}

\subsubsection{Proveedor}
Emisores de facturas y comprobantes.
\begin{itemize}
  \campo{id_proveedor}{Identificador único}{SERIAL, clave primaria}
  \campo{nombre}{Razón social o nombre}{VARCHAR(150)}
  \campo{tipo_doc_fiscal}{NIT o CI}{VARCHAR(3), CHECK}
  \campo{numero_doc_fiscal}{Número del documento fiscal}{VARCHAR(50)}
\end{itemize}

\subsubsection{Gasto}
Cada gasto realizado durante un viaje. Guarda los valores de retención calculados al registrarlo.
\begin{itemize}
  \campo{id_gasto}{Identificador único}{SERIAL, clave primaria}
  \campo{monto_total}{Importe total}{DECIMAL(10,2)}
  \campo{fecha_gasto}{Fecha del gasto}{DATE}
  \campo{descripcion}{Descripción}{VARCHAR(1000)}
  \campo{tipo}{F, R, C o S}{CHAR(1), CHECK}
  \campo{modificado}{Los datos extraídos se corrigieron a mano}{BOOLEAN}
  \campo{es_gasto_internacional}{Gasto en dólares}{BOOLEAN}
  \campo{moneda}{BOB o la moneda extranjera}{VARCHAR(10)}
  \campo{tipo_cambio}{Tipo de cambio aplicado}{NUMERIC(10,4)}
  \campo{monto_moneda_origen}{Importe en la moneda de origen}{NUMERIC(10,2)}
  \campo{base_imponible}{Base de cálculo de las retenciones}{NUMERIC(10,2)}
  \campo{retencion_rc_iva}{Retención RC-IVA}{NUMERIC(10,2)}
  \campo{retencion_iue}{Retención IUE}{NUMERIC(10,2)}
  \campo{retencion_it}{Retención IT}{NUMERIC(10,2)}
  \campo{importe_costo}{Importe imputable como costo}{NUMERIC(10,2)}
  \campo{tiene_alcohol}{Contiene bebidas alcohólicas}{BOOLEAN}
  \campo{id_viaje}{Viaje al que pertenece}{INTEGER}
  \campo{id_categoria}{Categoría del gasto}{INTEGER}
  \campo{id_proveedor}{Proveedor, opcional}{INTEGER}
\end{itemize}

\subsubsection{Gasto\_Subitem}
Desglose de un gasto en varios conceptos. Cuando existen, el monto total del gasto es su suma.
\begin{itemize}
  \campo{id_subitem}{Identificador único}{SERIAL, clave primaria}
  \campo{id_gasto}{Gasto al que pertenece}{INTEGER, borrado en cascada}
  \campo{descripcion}{Concepto}{VARCHAR(255)}
  \campo{monto}{Importe}{NUMERIC(10,2)}
\end{itemize}

\subsubsection{Gasto\_Tramo\_Moneda}
Tramos de conversión de moneda de un gasto internacional, como referencia del cálculo. El tipo de cambio lo escribe el empleado en cada tramo.
\begin{itemize}
  \campo{id_tramo}{Identificador único}{SERIAL, clave primaria}
  \campo{id_gasto}{Gasto al que pertenece}{INTEGER, borrado en cascada}
  \campo{moneda}{Moneda de origen}{VARCHAR(10)}
  \campo{monto_origen}{Importe en la moneda de origen}{NUMERIC(12,2)}
  \campo{tipo_cambio}{Tipo de cambio aplicado (unidades por 1 dólar)}{NUMERIC(12,4)}
  \campo{monto_usd}{Equivalente en dólares}{NUMERIC(12,2)}
\end{itemize}

\subsubsection{Imagen}
Comprobante adjunto a un gasto.
\begin{itemize}
  \campo{id_imagen}{Identificador único}{SERIAL, clave primaria}
  \campo{url_archivo}{URL del archivo}{TEXT}
  \campo{id_gasto}{Gasto al que pertenece}{INTEGER, único}
\end{itemize}

\subsubsection{Factura}
Factura vinculada a un gasto, en relación uno a uno.
\begin{itemize}
  \campo{id_factura}{Identificador único}{SERIAL, clave primaria}
  \campo{numero_factura}{Número del comprobante o, en recibos y boletos sin número, su código (admite símbolos)}{VARCHAR(50)}
  \campo{fecha_emision}{Fecha de emisión}{DATE}
  \campo{monto_parcial}{Importe sin impuestos}{DECIMAL(10,2)}
  \campo{id_gasto}{Gasto asociado}{INTEGER, único}
  \campo{id_proveedor}{Proveedor emisor}{INTEGER}
  \item Unicidad sobre \code{(numero_factura, fecha_emision, id_proveedor)} para evitar duplicados.
\end{itemize}

\subsubsection{Detalle\_Factura}
Productos o servicios incluidos en una factura.
\begin{itemize}
  \campo{id_detalle}{Identificador único}{SERIAL, clave primaria}
  \campo{nombre_producto}{Nombre del ítem}{VARCHAR(255)}
  \campo{cantidad}{Cantidad, con hasta tres decimales (litros, kilos o unidades)}{NUMERIC(12,3)}
  \campo{precio}{Precio unitario}{DECIMAL(10,2)}
  \campo{id_factura}{Factura a la que pertenece}{INTEGER}
\end{itemize}

\subsubsection{Impuesto y Factura\_Impuestos}
Catálogo de impuestos aplicables (\code{id_impuesto}, \code{nombre}, \code{porcentaje}) y tabla pivote que los relaciona con las facturas, con clave primaria compuesta \code{(id_factura, id_impuesto)}.

\subsubsection{Comentario}
Observaciones de las revisiones y justificaciones de exceso.
\begin{itemize}
  \campo{id_comentario}{Identificador único}{SERIAL, clave primaria}
  \campo{descripcion}{Contenido}{VARCHAR(300)}
  \campo{fecha}{Fecha de registro}{TIMESTAMPTZ}
  \campo{tipo}{JUSTIFICACION u OBSERVACION}{VARCHAR(20)}
  \campo{ciclo_revision}{Ronda en que se emitió}{INTEGER}
  \campo{fecha_justificada}{Día justificado; vacío para el exceso en hoteles}{DATE}
  \campo{id_usuario}{Autor}{INTEGER}
  \campo{id_viaje}{Viaje al que se refiere}{INTEGER}
  \campo{id_gasto}{Gasto observado; si se borra el gasto, la observación se conserva}{INTEGER}
\end{itemize}

\subsubsection{Codigo\_Verificacion}
Códigos de un solo uso para la recuperación de cuenta, válidos por cinco minutos.
\begin{itemize}
  \campo{id_codigo}{Identificador único}{SERIAL, clave primaria}
  \campo{codigo}{Código}{VARCHAR(7)}
  \campo{expiracion}{Vencimiento}{TIMESTAMPTZ}
  \campo{activo}{Indica si sigue siendo válido}{BOOLEAN}
  \campo{id_usuario}{Usuario al que pertenece}{INTEGER}
\end{itemize}

\subsubsection{Auditoria}
Eventos de acceso y de seguridad, con fines de trazabilidad.
\begin{itemize}
  \campo{id_auditoria}{Identificador único}{SERIAL, clave primaria}
  \campo{fecha}{Fecha del evento}{TIMESTAMPTZ}
  \campo{tipo}{INGRESO, SALIDA o CAMBIO\_CLAVE}{VARCHAR(15)}
  \campo{id_usuario}{Usuario que generó el evento}{INTEGER}
\end{itemize}

\subsubsection{Correlativo\_Recibo y Recibo}
\code{Correlativo_Recibo} es una tabla de una sola fila con el último número de recibo emitido. Su incremento lo hace una función almacenada que actualiza y devuelve el nuevo valor en una sola operación, evitando que dos emisiones simultáneas obtengan el mismo número. \code{Recibo} guarda el número asignado a cada recibo individual o agrupado, para reutilizarlo cuando se vuelve a enviar.
\begin{itemize}
  \campo{id_recibo}{Identificador único}{SERIAL, clave primaria}
  \campo{numero}{Número del recibo}{INTEGER, único}
  \campo{id_viaje}{Viaje}{INTEGER}
  \campo{id_gasto}{Gasto, vacío en los recibos agrupados}{INTEGER}
  \campo{tipo}{C (compra) o S (servicio)}{CHAR(1)}
  \campo{es_gasto_internacional}{Recibo de gastos en dólares}{BOOLEAN}
  \campo{fecha_emision}{Fecha de emisión}{TIMESTAMPTZ}
\end{itemize}

\subsubsection{Configuracion\_Recordatorio}
Tabla de una sola fila con los días y horas en que se envía el resumen de pendientes. El administrador la modifica desde la pantalla \emph{Recordatorios} y el servidor reprograma los envíos al guardar, sin reiniciarse. Si la fila no existe, se usa la configuración por defecto (lunes a viernes, 08:00, 12:00 y 16:00).
\begin{itemize}
  \campo{id}{Siempre 1}{SMALLINT, clave primaria}
  \campo{activo}{Indica si los envíos automáticos están activos}{BOOLEAN}
  \campo{dias}{Días de la semana (0 domingo a 6 sábado)}{SMALLINT[]}
  \campo{horas}{Horas de envío en formato HH:MM, hasta cuatro}{TEXT[]}
  \campo{fecha_actualizacion}{Fecha del último cambio}{TIMESTAMPTZ}
  \campo{id_usuario_actualizacion}{Administrador que hizo el último cambio}{INTEGER}
\end{itemize}

\subsection{Índices}
Se definieron índices sobre las columnas más usadas para filtrar o unir: las claves foráneas de las tablas dependientes, el estado del viaje, el supervisor asignado, el jefe directo y la sección del usuario, el estado de las solicitudes y el historial de revisión por persona y etapa.

\subsection{Seguridad a nivel de fila}
Se habilitó \textbf{Row Level Security} en las 24 tablas. El backend accede con la clave de servicio, que omite estas políticas; habilitarlas garantiza que ninguna conexión externa pueda leer o modificar datos directamente, y deja al servidor de aplicaciones como único punto de acceso.

\subsection{Scripts de base de datos}
El repositorio incluye los scripts para reconstruir el esquema desde cero, que se ejecutan en este orden:

\begin{xltabular}{\textwidth}{L{5.4cm} Y}
\cabecera \tth{Archivo} & \tth{Contenido} \\
\endfirsthead
\cabecera \tth{Archivo} & \tth{Contenido} \\
\endhead
\rowcolor{maxamFondo} \code{01_schema.sql} & Tablas, restricciones, claves foráneas e índices. \\
\code{02_functions.sql} & Función de incremento del correlativo de recibos. \\
\rowcolor{maxamFondo} \code{03_seed.sql} & Datos iniciales: roles, cargos, categorías, impuestos y configuración de recordatorios. \\
\code{04_rls_hardening.sql} & Activación de la seguridad a nivel de fila. \\
\rowcolor{maxamFondo} \code{testing_*.sql} & Limpieza y organización de prueba; solo para entornos de prueba y para las pruebas automáticas. \\
\end{xltabular}

\subsection{Acceso a la base de datos}
La base de datos está alojada en \textbf{Supabase}, plataforma de backend como servicio construida sobre PostgreSQL que ofrece administración visual de tablas, autenticación, almacenamiento y acceso mediante API REST. El panel de administración del proyecto está disponible en:

\begin{center}
\url{https://supabase.com/dashboard/project/ubxqgumxpflxzsigdcxj}
\end{center}

Las credenciales de la base de datos y las claves de los servicios externos se gestionan mediante variables de entorno y no se incluyen en esta documentación.

\end{document}
